\chapter{Tests et validation du système}

\section{Stratégie de qualification et plan de tests}

La phase d'assurance qualité (\emph{QA}) démontre la conformité de la plateforme vis-à-vis des spécifications fonctionnelles (\nolinkurl{docs/requirements.md}) et des exigences d'étanchéité multi-établissements. Conformément à la déontologie académique, aucun résultat de test hypothétique n'est consigné. Chaque métrique présentée émane directement du journal d'exécution horodaté (\nolinkurl{logs/qa_execution.log}) et des preuves machine (\nolinkurl{docs/evidence/}).

La stratégie de qualification repose sur une pyramide de tests à quatre niveaux :
\begin{enumerate}
    \item \textbf{Tests d'intégration HTTP automatisés} : 16 suites de tests exécutant des requêtes REST réelles et validant le code de retour, le format JSON et les en-têtes via la bibliothèque d'assertions \texttt{lib\_http\_assert.sh}.
    \item \textbf{Tests d'étanchéité multi-tenant (SEC-003)} : scénarios spécifiques confrontant deux établissements distincts (\texttt{ecole-221} et \texttt{ecole-test-b}) pour certifier l'impossibilité d'une fuite de données inter-écoles.
    \item \textbf{Contrôles statiques et règles d'agents} : vérification automatisée de l'absence totale de l'en-tête \texttt{X-Tenant-Id} côté frontend et respect des cibles de typage sans \texttt{any} (\texttt{audit/check\_agent\_rules.py}).
    \item \textbf{Validation navigateur de bout en bout (E2E)} : automatisation de parcours réels sous Playwright (\nolinkurl{qa/e2e/parcours_navigateur.py}) simulant la navigation d'un administrateur, d'un élève et d'un vigile.
\end{enumerate}

\section{Environnement et jeux de données certifiés}

Pour éviter toute pollution des données de production ou incohérence entre exécutions, l'ensemble des suites s'exécute dans un environnement dédié :
\begin{itemize}
    \item \textbf{Serveur d'application} : instance Spring Boot compilée et démarrée sur le port local 8090, configurée avec une base de données PostgreSQL isolée (port 55433).
    \item \textbf{Jeu de données certifié (\texttt{qa/fixtures.sql})} : script d'initialisation reproductible provisionnant deux écoles cloisonnées (\texttt{ecole-221} et \texttt{ecole-test-b}), les comptes administrateurs de test, les filières, les enseignants et les élèves modèles.
    \item \textbf{Authentification sécurisée} : les mots de passe de test sont rigoureusement synchronisés via des variables d'environnement locales, éliminant tout secret en dur dans le code source.
\end{itemize}

\section{Résultats expérimentaux des tests d'intégration}

L'exécution intégrale de la suite de qualification via la commande \texttt{./orchestrator.sh test} aboutit à un taux de succès de \textbf{100\,\%}, validé par le verdict officiel \texttt{QA\_RESULT: PASS} consigné dans \nolinkurl{docs/test-report.md}.

Le tableau~\ref{tab:resultats-qa} détaille les résultats obtenus pour chacune des 16 suites de tests du progiciel.

\begin{table}[h]
\centering
\small
\begin{tabular}{llp{5.5cm}cc}
\toprule
\textbf{Suite} & \textbf{Script de test} & \textbf{Périmètre fonctionnel couvert} & \textbf{Assertions} & \textbf{Résultat} \\
\midrule
QA-01 & \texttt{test\_auth\_module.sh} & Authentification JWT, cookies \& inscription & 9 / 9 & \textbf{PASS} \\
QA-02 & \texttt{test\_superadmin\_tenants.sh} & Gestion des écoles \& quotas de plateforme & 8 / 8 & \textbf{PASS} \\
QA-03 & \texttt{test\_billing\_module.sh} & Souscriptions, formules SaaS \& abonnements & 7 / 7 & \textbf{PASS} \\
QA-04 & \texttt{test\_admin\_module.sh} & Administration scolaire, filières \& classes & 10 / 10 & \textbf{PASS} \\
QA-05 & \texttt{test\_student\_module.sh} & Espace étudiant, notes \& tableau de bord & 6 / 6 & \textbf{PASS} \\
QA-06 & \texttt{test\_professor\_module.sh} & Cours, devoirs \& saisie des notes CC/examen & 8 / 8 & \textbf{PASS} \\
QA-07 & \texttt{test\_parent\_module.sh} & Suivi tuteur, présences \& paiements & 7 / 7 & \textbf{PASS} \\
QA-08 & \texttt{test\_security\_module.sh} & Postes vigiles, scans QR badges \& alertes & 7 / 7 & \textbf{PASS} \\
QA-09 & \texttt{test\_finance\_module.sh} & Dépenses, encaissements \& trésorerie & 7 / 7 & \textbf{PASS} \\
QA-10 & \texttt{test\_planning\_module.sh} & Emplois du temps \& gestion des conflits & 7 / 7 & \textbf{PASS} \\
QA-11 & \texttt{test\_candidatures\_module.sh} & Admissions en ligne \& conversion en élève & 18 / 18 & \textbf{PASS} \\
QA-12 & \texttt{test\_collaboration\_module.sh} & Salons de messagerie \& documents partagés & 8 / 8 & \textbf{PASS} \\
QA-13 & \texttt{test\_dashboard\_module.sh} & Consolidation des indicateurs clés (KPIs) & 6 / 6 & \textbf{PASS} \\
QA-14 & \texttt{test\_final\_modules.sh} & Cohérence croisée \& intégrité référentielle & 10 / 10 & \textbf{PASS} \\
QA-15 & \texttt{test\_referentiels\_module.sh} & Référentiels d'école \& listes paramétrables & 28 / 28 & \textbf{PASS} \\
QA-16 & \texttt{test\_tenant\_isolation\_module.sh} & Étanchéité stricte multi-tenant A vs B & 9 / 9 & \textbf{PASS} \\
\midrule
\textbf{Total} & \textbf{16 modules certifiés} & \textbf{Couverture fonctionnelle intégrale} & \textbf{142 / 142} & \textbf{100\,\% PASS} \\
\bottomrule
\end{tabular}
\caption{Matrice d'exécution des tests d'intégration certifiés}
\label{tab:resultats-qa}
\end{table}

\subsection{Démonstration de l'étanchéité multi-tenant (Module QA-16)}

Le module QA-16 revêt une importance capitale : il simule une tentative d'accès frauduleux où un administrateur authentifié auprès de l'école A tente d'interroger ou de modifier un dossier d'étudiant appartenant à l'école B. Les 9 scénarios du script vérifient que :
\begin{enumerate}
    \item Une tentative d'accès direct par identifiant technique est rejetée avec un code HTTP \texttt{403 Forbidden}.
    \item Les requêtes de liste (\texttt{GET /api/student}) exécutées avec le token de l'école A ne renvoient strictement aucun résultat rattaché à l'école B.
    \item L'injection d'un en-tête malveillant \texttt{X-Tenant-Slug} non concordant avec le token déclenche le blocage immédiat par le \texttt{TenantMembershipInterceptor}.
\end{enumerate}

\section{Anomalies, corrections et tests de non-régression}

Le cycle d'assurance qualité a mis en lumière trois anomalies critiques qui ont été résolues et vérifiées :
\begin{itemize}
    \item \textbf{Fermeture de l'accès public \nolinkurl{/admissions/db} (APP-TL-10)} : une route technique de diagnostic exposait la liste des dossiers d'admissions sans authentification préalable. L'interdiction d'accès public a été durcie dans \texttt{SecurityConfig.java} et validée par test unitaire.
    \item \textbf{Gestion du stockage Cloudinary en environnement déconnecté (APP-BE-05)} : en l'absence de connectivité Internet sur le banc de test local, le service de téléversement levait une erreur 503. Un mode de simulation déterministe (\texttt{CLOUDINARY\_SIMULATION=true}) a été intégré pour garantir la reproductibilité des tests sans dépendance externe.
    \item \textbf{Résorption de la dette technique Frontend} : élimination de 10 occurrences de types \texttt{any}, intégration de Zod et React Hook Form, et validation d'une compilation TypeScript (\texttt{tsc -b}) à 0 erreur.
\end{itemize}

\section{Bilan de conformité et limites de la validation}

Les 142 assertions exécutées avec succès démontrent que les objectifs spécifiques formulés dans l'introduction générale (authentification RBAC, autonomie d'inscription des écoles, intégrité des notes et étanchéité multi-tenant) sont pleinement atteints.

Toutefois, une démarche académique honnête impose de souligner les limites actuelles du banc d'essai :
\begin{itemize}
    \item Les tests de charge sous fort trafic simultané (plus de 10\,000 requêtes concurrentes) n'ont pas été menés en environnement réel, l'infrastructure de test étant hébergée sur une machine locale de développement.
    \item L'évaluation de l'expérience utilisateur et de l'accessibilité ergonomique a été menée sur un échantillon restreint d'utilisateurs pilotes, sans campagne d'évaluation standardisée à grande échelle.
\end{itemize}
